The efficiency case for sparsifying inter-agent attention rests on an operation
count. We examined that count and then measured what runs.

The count is smaller than reported: because selecting peers by attention score
requires every score, only half of attention is sparsifiable and the achievable
FLOPs reduction is bounded by $2\times$, not by $1 - k/(N-1)$. The count is
also, in the implementation pattern we found in use, not executed at all ---
masking the dense attention matrix leaves the dense aggregation matmul in
place. And when we implemented the sparse gather the claim actually refers to,
verified it computes the same function, and measured it against a fused dense
baseline, it was slower in every one of $46$ cells spanning head width, batch,
agent count and sparsity level, on two architectures, at two precisions, under
three execution modes. Selection alone costs an order of magnitude more than
the aggregation it is meant to shrink, the sparse contraction is itself slower
than the dense one, and both paths are memory-bound --- so removing arithmetic
was never going to help.

None of this shows that sparse multi-agent communication is a bad idea. It
shows that the standard argument for it does not survive measurement in the
centralised-inference regime, and it identifies precisely what a successor
method must do differently. Proposition~\ref{prop:escape} is the constructive
half of the result: a method whose sparsity pattern is fixed \emph{before} the
scores are computed --- by topology, by a spatial prior, by a hash --- is not
subject to the ceiling and can reach $N/k$. In an environment with spatial
structure, distance-calibrated routing is the natural instance, and it is where
we would spend the next unit of effort. A fused kernel that eliminates the
intermediate traffic between selection, gather and contraction is the other.

The wider lesson is methodological, and it is not specific to MARL. An
operation count is a model of cost. Like any model it can be wrong in two ways:
it can mis-estimate the quantity it describes, and the quantity it describes
can fail to govern the outcome anyone cares about. Both happened here. The
remedy is cheap --- report executed operation counts beside measured latency,
so that a divergence between the algorithm claimed and the algorithm executed
becomes visible immediately. Our audit suggests the surrounding literature has
largely not been doing this: one of six multi-agent communication papers
reports any measured temporal quantity for its mechanism, against five of six
in the efficient-attention work the argument was borrowed from. We release the
harness, the corrected implementation, the pre-registration and the five
benchmarking pitfalls that produced spurious results along the way.
